\begin{lemma}
\label{lemma-automatic-connection}
In Situation \ref{situation-global}.
Let $\mathcal{F}$ be a crystal in $\mathcal{O}_{X/S}$-modules
on $\text{Cris}(X/S)$. Then $\mathcal{F}$ comes equipped with a
canonical integrable connection.
\end{lemma}

\begin{proof}
Say $(U, T, \delta)$ is an object of $\text{Cris}(X/S)$.
Let $(U, T', \delta')$ be the infinitesimal thickening of $T$
by $(\Omega_{X/S})_T = \Omega_{T/S, \delta}$
constructed in Remark \ref{remark-first-order-thickening}.
It comes with projections $p_0, p_1 : T' \to T$
and a diagonal $i : T \to T'$. By assumption we get
isomorphisms
$$
p_0^*\mathcal{F}_T \xrightarrow{c_0}
\mathcal{F}_{T'} \xleftarrow{c_1}
p_1^*\mathcal{F}_T
$$
of $\mathcal{O}_{T'}$-modules. Pulling $c = c_1^{-1} \circ c_0$
back to $T$ by $i$ we obtain the identity map
of $\mathcal{F}_T$. Hence if $s \in \Gamma(T, \mathcal{F}_T)$
then $\nabla(s) = p_1^*s - c(p_0^*s)$ is a section of
$p_1^*\mathcal{F}_T$ which vanishes on pulling back by $i$. Hence
$\nabla(s)$ is a section of
$$
\mathcal{F}_T
\otimes_{\mathcal{O}_T}
\Omega_{T/S, \delta}
$$
because this is the kernel of $p_1^*\mathcal{F}_T \to \mathcal{F}_T$
as $\mathcal{O}_{T'} = \mathcal{O}_T \oplus \Omega_{T/S, \delta}$
by construction. It is easily verified that $\nabla(fs) =
f\nabla(s) + s \otimes \text{d}(f)$ using the description of
$\text{d}$ in Remark \ref{remark-first-order-thickening}.

\medskip\noindent
The collection of maps
$$
\nabla : \Gamma(T, \mathcal{F}_T) \to
\Gamma(T, \mathcal{F}_T \otimes_{\mathcal{O}_T} \Omega_{T/S, \delta})
$$
so obtained is functorial in $T$ because the construction of $T'$
is functorial in $T$. Hence we obtain a connection.

\medskip\noindent
To show that the connection is integrable we consider the
object $(U, T'', \delta'')$ constructed in
Remark \ref{remark-second-order-thickening}.
Because $\mathcal{F}$ is a sheaf we see that
$$
\xymatrix{
q_0^*\mathcal{F}_T \ar[rr]_{q_{01}^*c} \ar[rd]_{q_{02}^*c} & &
q_1^*\mathcal{F}_T \ar[ld]^{q_{12}^*c} \\
& q_2^*\mathcal{F}_T
}
$$
is a commutative diagram of $\mathcal{O}_{T''}$-modules. For
$s \in \Gamma(T, \mathcal{F}_T)$ we have
$c(p_0^*s) = p_1^*s - \nabla(s)$. Write
$\nabla(s) = \sum p_1^*s_i \cdot \omega_i$ where $s_i$ is a local section
of $\mathcal{F}_T$ and $\omega_i$ is a local section of $\Omega_{T/S, \delta}$.
We think of $\omega_i$ as a local section of the structure
sheaf of $\mathcal{O}_{T'}$ and hence we write product instead of tensor
product. On the one hand
\begin{align*}
q_{12}^*c \circ q_{01}^*c(q_0^*s) & = 
q_{12}^*c(q_1^*s - \sum q_1^*s_i \cdot q_{01}^*\omega_i) \\
& =
q_2^*s - \sum q_2^*s_i \cdot q_{12}^*\omega_i -
\sum q_2^*s_i \cdot q_{01}^*\omega_i +
\sum q_{12}^*\nabla(s_i) \cdot q_{01}^*\omega_i
\end{align*}
and on the other hand
$$
q_{02}^*c(q_0^*s) = q_2^*s - \sum q_2^*s_i \cdot q_{02}^*\omega_i.
$$
From the formulae of Remark \ref{remark-second-order-thickening} we see
that
$q_{01}^*\omega_i + q_{12}^*\omega_i - q_{02}^*\omega_i = \text{d}\omega_i$.
Hence the difference of the two expressions above is
$$
\sum q_2^*s_i \cdot \text{d}\omega_i -
\sum q_{12}^*\nabla(s_i) \cdot q_{01}^*\omega_i
$$
Note that
$q_{12}^*\omega \cdot q_{01}^*\omega' = \omega' \wedge \omega =
- \omega \wedge \omega'$ by the definition of the multiplication on
$\mathcal{O}_{T''}$. Thus the expression above is $\nabla^2(s)$ viewed
as a section of the subsheaf $\mathcal{F}_T \otimes \Omega^2_{T/S, \delta}$ of
$q_2^*\mathcal{F}$. Hence we get the integrability condition.
\end{proof}